% Post-CP25 theorem; exact audited candidate retained separately.
\subsection{A uniform nodal norm-seven self-isogeny in two-rank zero}
\label{subsec:uniform-seven-theta-candidate}
\status{Proved and independently checked.}

All geometric assertions in this subsection are over an algebraically
closed field of characteristic two. Write $\iota$ for the hyperelliptic
involution and use the unique Weierstrass point $W$ as Abel--Jacobi base
point. The two-rank-zero genus-two normal form and its presentation
\[
 C_z:\quad y^2+y=ax^5+bx^3,\qquad
 a=z^2,\quad b=z^2+z+1,\quad z\ne0
 \tag{S7.1}
\]
were proved in the uniform norm-five subsection. In particular a fixed
normalized modulus $\mu$ has the presentations determined by
\[
             (z^2+z+1)^5=\mu^5z^6.                \tag{S7.2}
\]
The change of the fifth root used to normalize the leading coefficient
permits the prescribed value of $\mu$, not merely its fifth power.

\begin{lemma}[Quaternion parameters and exact interpolation]
\label{lem:uniform-seven-interpolation-candidate}
\textup{[Proved; independent audits passed.]}
Choose $r$ with $zr^3+r^2+1=0$, put $D=r^2$, and choose
\[
 t^2+t=D,\qquad \beta^2+\beta=z+1,\qquad
 w^2+w=at^5+bt^3.
\]
Set
\[
 G(x)=zx^2+(z+1)x+\beta,\qquad
 H(x)=ux^2+vx+w,
 \quad u=z(t+r),\quad v=zt^2+t+r^{-2}.
\]
The automorphisms $g:(x,y)\mapsto(x+1,y+G(x))$ and
$h:(x,y)\mapsto(x+t,y+H(x))$ satisfy
$g^2=h^2=\iota$ and $gh=\iota hg$.
Thus $\psi=2+g+h+gh$ is an etale polarized self-isogeny of
$J(C_z)$, of multiplier seven and degree $49$.

Suppose additionally $c_g=zD+t\ne0$, and put $c=Dc_g$.
For a finite point $P=(p,y_0)$ define
\[
\begin{gathered}
 X=x+p,\qquad d_1=ap^4+bp^2,\\
 Z=H(p)+t^2G(p)+Dd_1+tc_g,\qquad
 E=tG(p)+H(p)+(t+1)c_g,\\
 A_3=(ZE+cc_g)/D,\quad
 A_2=(Z+c)G(p)+cd_1+A_3,\quad A_1=cd_1.
\end{gathered}
\]
Then $f_P=(ZX+c)(y+y_0)+A_1X+A_2X^2+A_3X^3$ has divisor
$D_P+E_P-7W$, where $D_P=2P+gP+hP+ghP$ and $E_P$ is an
effective nonspecial divisor of degree two. Its norm is
\[
 f_P\iota(f_P)=X^2(X+1)(X+t)(X+t+1)
 \bigl(aZ^2X^2+(A_3^2+apZ^2)X
              +(b+a(D+1))Z^2+ac^2\bigr).          \tag{S7.3}
\]
Moreover $\ker\psi\cap\Theta=\{W\}$.
\end{lemma}

\begin{proof}
The cubic equation for $r$ implies $r\ne0,1$ and $z=(r^2+1)/r^3$.
The exact identities
\[
 u^2=at,\quad v^2+u=bt,\quad v=at^4+bt^2,
 \quad ut^2+vt=1,\quad zt^2+(z+1)t+u+v=1
\]
give the curve equations, squares and commutator. Their actions on the
Jacobian satisfy the quaternion relations with $\iota=[-1]$.
Automorphisms preserving $W$ preserve the principal polarization, so
Rosati sends $g,h,gh$ to their negatives. Direct multiplication gives
$\psi^\dagger\psi=[7]$, proving the isogeny and degree statements.

The basis $1,x,x^2,x^3,y,xy$ of $H^0(7W)$ shows directly that an
arbitrary section vanishing on $D_P$ has the form
\[
 (LX+C)(y+y_0)+Cd_1X+B_2X^2+B_3X^3,
\]
with its three remaining value conditions equivalent to
\[
 CZ=LDc_g,\qquad
 B_3=(LE+Cc_g)/D,\qquad
 B_2=(L+C)G(p)+Cd_1+B_3.                         \tag{S7.4}
\]
For an explicit elimination, first use the value at $X=1$ to obtain
the displayed formula for $B_2$. If $E_h,E_{gh}$ are the two
remaining values at $X=t,t+1$, substitution gives
\[
 E_h=t(DB_3+LE+Cc_g)+(CZ+LDc_g),\qquad
 E_h/t+E_{gh}/(t+1)=(CZ+LDc_g)/D.
\]
All three denominators $t,t+1,D$ are nonzero, so these are reversible
row operations and prove the asserted equivalence.
The four prescribed $x$-coordinates $p,p+1,p+t,p+t+1$ are distinct,
because $D\ne0$. Since $Dc_g\ne0$, (S7.4) has a one-dimensional
solution space, including when $Z=0$. It is spanned by $f_P$.
Consequently $h^0(7W-D_P)=1$, so $E_P$ is nonspecial.
Expanding the norm gives (S7.3); the retained universal interpolation
certificate checks all these identities, including every prescribed
value and the derivative at $P$.

The class of $E_P-2W$ is $-\psi(P-W)$. If a finite $P$ were in
$\ker\psi\cap\Theta$, then $D_P\sim5W$, which would imply
$h^0(7W-D_P)=h^0(2W)=2$, a contradiction.
\end{proof}

\begin{lemma}[A necessary polynomial obstruction to a triple fibre]
\label{lem:uniform-seven-triple-obstruction-candidate}
\textup{[Proved; independent audits passed.]}
Use the preceding parameters with $c_g\ne0$. Put
\[
\begin{gathered}
 Z_0=w+t^2\beta+tc_g,\quad
 E_0=w+t\beta+(t+1)c_g,\quad q_0=(E_0^2+Z_0)/D^2,\\
 A_4=aD,\quad A_2^*=aD(D+1)+1,\quad A_1^*=aD^2,
 \quad A(L)=A_4L^3+A_2^*L+A_1^*,\\
 j=aD^2L+D^2q_0,\quad k=D^2c_g^4,\quad
 H=c_g^2(c_g^2+A_4L^2+A_2^*),\\
 P_2=c_g^4(c_g^2+A_4L^2+A_2^*)^2+A_4^2j^2
                         +LA(L)A_4c_g^4,\\
 P_1=A_4^2k+c_g^6A(L)(Lc_g^2+A_1^*),\qquad
 T=H(H+A_4j)^2+kA_4^3.
\end{gathered}                                                   \tag{S7.5}
\]
A triple fibre of $\psi|_\Theta$ implies either that $P_1,P_2$
have a common root $L$, or that $A,T$ have a common root $L$.
\end{lemma}

\begin{proof}
First dispose of the omitted coordinates. Infinity $W$ has no other
preimage by the preceding lemma. Two distinct finite points with the
same $x$-coordinate are hyperelliptic conjugates. If their images
coincided, that image would be geometric two-torsion; the two-rank is
zero, so it would be zero, again impossible.

The interpolation identities give
\[
 Z=A_4p^4+A_2^*p^2+A_1^*p+Z_0,
 \qquad E^2+apD^2=Z+E_0^2+Z_0.                    \tag{S7.6}
\]
Thus $Z$ has derivative $A_1^*\ne0$ and four distinct roots. At a
root of $Z$, one has $A_3=c_g^2\ne0$: the function has pole order
six, and $E_P=W+R$ with
$x(R)=p+aD^2/c_g^2$. These residual divisors have pairwise distinct
$x(R)$ as $p$ varies, and the two points above a fixed $p$ have
distinct images as just proved. None agrees with a residual divisor
at $Z\ne0$, which is finite of degree two. The uniqueness used here
is exactly the nonspeciality already proved, including for a doubled
effective residual point.

It remains to consider $Z\ne0$. If several points have a common
image, their unique residual divisors have the same monic quadratic
$x^2+Lx+M$. Set $\kappa=(b+a(D+1))/a$ and
$\delta=M+\kappa$. The two coefficients of (S7.3), in the global
$x$-coordinate, give
\[
 L=\frac{Z}{aD^2}+\frac{q_0}{a}+\frac{c_g^4}{aZ^2},\qquad
 p^2+Lp=\delta+\frac{c^2}{Z^2}.                  \tag{S7.7}
\]
Consequently every such $Z$ is a root of
$C(Z)=Z^3+jZ^2+k$. Its constant term $k$ is nonzero and its
derivative is $Z^2$, so its three roots are distinct and nonzero.
In its quotient algebra,
$Z^{-2}=(Z+j)/k$. Write
\[
 e=c_g^{-2},\qquad u_0=\delta+ej,
 \qquad u(Z)=u_0+eZ.
\]
Substitution of $p^2+Lp=u(Z)$ into (S7.6) gives
\[
 A(L)p=B(Z)=B_0+B_1Z+B_2Z^2,                    \tag{S7.8}
\]
where
\[
 B_0=Z_0+(A_4L^2+A_2^*)u_0+A_4u_0^2,\quad
 B_1=1+e(A_4L^2+A_2^*),\quad B_2=A_4e^2\ne0.
\]
If $A(L)\ne0$, each $Z$ supplies at most one $p$. A triple
therefore uses all three roots of $C$. The polynomial
$B^2+LA(L)B+uA(L)^2$, reduced modulo $C$, vanishes at all three
roots and hence is zero. Its $Z^2$ and $Z$ coefficients, multiplied
by $c_g^8$, are precisely $P_2$ and $P_1$. This includes $L=0$;
no division by $L$ occurs.

If $A(L)=0$, every possible $Z$ is a root of the genuine quadratic
$B$, and each such $Z$ supplies at most two values of $p$ by (S7.7).
A triple must use two distinct roots of $B$, so $B/B_2$ divides $C$.
Write $B/B_2=Z^2+hZ+i$, with $h=H/A_4$. Polynomial division gives
$i=h^2+jh$ and $h(h+j)^2+k=0$. Clearing the nonzero denominator
$A_4^3$ gives $T=0$. This also handles a repeated root of $B$:
in that case there is only one $Z$ and at most two points, so it
cannot give a triple.
\end{proof}

\begin{lemma}[Exact universal parameter certificate]
\label{lem:uniform-seven-parameter-certificate-candidate}
\textup{[Proved; independent audits passed.]}
There is a squarefree polynomial $\mathcal B(r)\in\mathbf F_2[r]$
of degree $97$, given by the retained universal certificate, with
the following properties. It is divisible by $r+1$ and
$r^3+r+1$. For $r\ne0,1$ with $\mathcal B(r)\ne0$, both roots
of $t^2+t=r^2$ have $c_g\ne0$, neither pair $P_1,P_2$ nor $A,T$
has a common root, and the theta image is therefore triple-free.
For $z\ne0$, all three roots of $zr^3+r^2+1$ can lie in
$\{\mathcal B=0\}$ only if
\[
                         z^3+z^2+1=0.           \tag{S7.9}
\]
\end{lemma}

\begin{proof}
Here are the operations and exact outputs needed to reproduce the
certificate; they involve universal polynomial identities, not a
sample of parameter values. Work in
$\mathbf F_2(r)[L,t]/(t^2+t+r^2)$. In this algebra,
\[
 q_0=\frac{r^2+1}{r^4}t+\frac{r^2+1}{r^6}.
\]
Pseudo-divide the degree-four $P_1$ by the degree-two $P_2$, with
no inversion of the leading coefficient of $P_2$, obtaining
$R_1L+R_0$. If $c_2L^2+c_1L+c_0=P_2$, form
\[
 O=c_2R_0^2+c_1R_0R_1+c_0R_1^2.
\]
A common root of $P_1,P_2$ forces $O=0$, even at a specialization
where a leading coefficient in the pseudo-division vanishes.
Separately form the determinant of multiplication by $T$ in the
rank-three algebra defined by the monic polynomial $A/A_4$.
A common root of $A,T$ forces that determinant to vanish. Take the
norms of these two quantities under $t\mapsto t+1$.
Their denominators are respectively $r^{172}$ and $r^{72}$.
Their numerator radicals have union
\[
\begin{aligned}
\mathcal B(r)={}&(r+1)(r^3+r+1)(r^6+r+1)(r^6+r^3+1)\\
 &\cdot(r^6+r^4+r^3+r+1)(r^9+r^4+r^2+r+1)\\
 &\cdot(r^{12}+r^9+r^8+r^7+r^6+r^4+1)\\
 &\cdot(r^{15}+r^{12}+r^{11}+r^9+r^7+r^6+r^3+r+1)\\
 &\cdot(r^{39}+r^{36}+r^{35}+r^{33}+r^{31}+r^{29}
       +r^{27}+r^{25}+r^{23}+r^{21}\\
 &\hspace{19mm}+r^{19}+r^{18}+r^{17}+r^{14}+r^{12}
       +r^8+r^7+r^6+r^3+r+1).
\end{aligned}                                                     \tag{S7.10}
\]
The norm of $c_g=t+zD$ is $(r^3+r+1)/r^2$, already included.
All divisions used above have denominators supported on $r=0,1$.
Thus nonvanishing of $\mathcal B$ gives all asserted exclusions.

Finally divide $\mathcal B(r)$ by the monic cubic
$r^3+z^{-1}r^2+z^{-1}$ over $\mathbf F_2(z)$. The greatest common
divisor of the three remainder-coefficient numerators is exactly
$z^3+z^2+1$. The coefficient denominators are powers of $z$.
The cubic is separable at every $z\ne0$. If all its roots are
roots of $\mathcal B$, the remainder is identically zero after
specialization, which proves (S7.9).

The complete arithmetic is retained in
\path{code/monodromy/verify_nonordinary_seven_good_parameters.py}, with exact
output, the two separate norm factorizations, the full remainder,
and checks of the displayed radical and gcd. The preceding
interpolation identities are independently presented in
\path{code/monodromy/nonordinary_seven_q8_interpolation.py}. These are the
only numerical-algebra inputs in this proof; all exceptional
geometric cases were dealt with before the certificate.
\end{proof}

\begin{theorem}[A nodal seven-self-isogeny for every two-rank-zero Jacobian]
\label{thm:uniform-seven-nodal-self-isogeny-candidate}
\textup{[Proved; independent audits passed.]}
Every indecomposable principally polarized abelian surface of
two-rank zero in characteristic two admits an etale polarized
self-isogeny of multiplier seven whose principal theta image is
integral nodal, with exactly $48$ nodes.
\end{theorem}

\begin{proof}
For $\mu\ne0$, the polynomial in (S7.2) has degree ten, nonzero
constant term, and derivative $(z^2+z+1)^4$, which is nonzero at
each of its roots. It therefore gives ten distinct nonzero
presentations, at least one of which avoids the three roots of
(S7.9). For $\mu=0$ the two presentations are the roots of
$z^2+z+1$, and neither satisfies (S7.9). Choose such a $z$.
The parameter certificate provides a root $r$ with
$\mathcal B(r)\ne0$. Either choice of $t$, and any choices of
$\beta,w$, then gives the stated quaternion self-isogeny.
The triple-fibre obstruction and certificate prove that every
theta fibre has at most two points. The proved odd-kernel theta
nodality criterion, Lemma~\ref{lem:root-nonordinary-odd-theta-nodality},
now gives an integral nodal image, with
$|\ker\psi|-1=49-1=48$ nodes.
\end{proof}

\begin{corollary}[The three--five--seven degree family]
\label{cor:f0-three-five-seven-degree-groups-candidate}
\textup{[Proved; independent audits passed.]}
For every fixed two-rank-zero principally polarized abelian surface
in characteristic two, the full geometric and arithmetic
Hasse--Witt monodromy conclusions hold for odd degrees
$d=3^i5^j7^k$, $i,j,k\ge0$, and for every even degree which is a
sum of two such odd degrees. In particular the previously open
fixed-surface degrees seven and twenty-two are included.
\end{corollary}

\begin{proof}
For indecomposable surfaces, iterate the proved good-theta norm
transfer using the nodal self-isogenies of multipliers five and
seven and the already established powers-of-three theorem.
The even conclusion is exactly the proved even-addition theorem.
For decomposable surfaces the all-degree product theorem already
gives the conclusion. The groups are $\operatorname{GL}_{9d^2-1}
(\mathbf F_2)$ for odd $d$ and
$\operatorname{Hom}(\mathbf F_2,\mathbf F_2^{9d^2-4})\rtimes
\operatorname{GL}_{9d^2-4}(\mathbf F_2)$ for even $d$.
The statements descend to any perfect ground field by the same
geometric/arithmetic bound and invariant-covector descent used
in the good-theta transfer theorem.
\end{proof}

\begin{remark}[Scope and retained explorations]
\label{rem:uniform-seven-scope-candidate}
\textup{[Proved; independent audits passed.]}
This degree family does not comprise all positive integers.
No statement for every odd kernel is made. Earlier fixed-$C_0$
Mumford-coordinate and Groebner calculations, including all failed
or timed-out versions, remain retained as exploratory backup;
they are not inputs to the universal proof above. The general
automorphism context in Brock--Howe, Lemma~2.1, pp.~4--8, was
actually read, but neither its computer-assisted census nor any
odd-norm nodality assertion is imported here. The quaternion and
parameter identities used here are verified directly. The complete parent and obstruction/crystal proof/code audits passed before promotion; both universal programs were rerun with the author outputs preserved. The author parameter certificate passed at 21:49:14 UTC in $0.18161791599413846$ seconds; the companion runs took $0.198577982$ and $0.306713550$ seconds. Exact records and all exploratory failures are retained in the agent output and history directories.
\end{remark}
